\section{Results}\label{sec:results}

\subsection{On holiday the network splits into senders and receivers}\label{sec:res-structure}

The summer network is not just the ordinary one with more people. On an ordinary weekday, a zone whose residents sleep in many places also receives people from many places: the rank correlation between out-degree and in-degree is 0.69, and the points of Figure~\ref{fig:senders}a lie along the diagonal. On 14--15 August the correlation falls to 0.09. The cities move below the diagonal: residents of Madrid sleep in 2,203 zones, but people from only 995 zones sleep in Madrid. Coastal towns move above it, and they make up 40 of the 50 zones that receive people from most places, against 18 in February. The typical zone now receives people from 27 zones and sends them to 14 (Table~\ref{tab:basic}). Reciprocity, the chance that a link $i\to j$ has a partner $j\to i$, halves from 0.52 to 0.24. This is not an effect of density: on the summer network cut to the same number of links as February, reciprocity is still 0.24 (Table~\ref{tab:basic_full}).

\begin{figure}[t]
\centering
\includegraphics[width=\linewidth]{fig2_senders_receivers.pdf}
\caption{\textbf{On the summer night the network splits into senders and receivers.} \textbf{a, b}, Out-degree against in-degree of each zone on the ordinary night (a) and on 14--15 August (b), on logarithmic axes; points are slightly jittered. \textbf{c}, Degree assortativity $r$ of the real network (dot), the mean of 50 degree-preserving rewirings (tick) and the range reached by maximally assortative and disassortative rewiring that keeps the network connected (bar; \citealp{zhou2007structural}).}
\label{fig:senders}
\end{figure}

The mixing of the network changes too. Degree assortativity falls from $-0.131$ on the ordinary night to $-0.214$ in summer: on holiday, the hubs link even more to small zones. But the new degrees force most of this. The range of $r$ that the degrees allow moves from $[-0.22,-0.09]$ to $[-0.42,-0.11]$, and random networks with the same degrees give $-0.140$ and $-0.217$, close to the real values (Figure~\ref{fig:senders}c). Rewiring reproduces about 93\,\% of the change from February. In the language of \citet{zhou2007structural}, the degree sequence constrains the mixing; the remaining gap, about six standard deviations of the rewired networks, is real but small.

The rich club tells a similar story. The 50 best-connected zones, mostly cities, are linked in 80\,\% of their possible pairs on the ordinary night and 96\,\% in summer. Random networks with the same degrees give 66\,\% and 92\,\%. So on an ordinary night the cities form a club beyond what their degrees produce; in summer, when they link to almost every zone, their degrees alone almost make them a clique. The bitmaps of the sorted adjacency matrix confirm the picture: the three nights have the same overall shape, the summer one denser (Figure~\ref{fig:bosam}).

One property does go beyond the degrees. Second-order mixing asks whether the neighbours of two linked zones are themselves alike. In summer it is $-0.289$, against $-0.252$ in rewired networks that keep both the degrees and the rich club; on the summer network with February's number of links it is $-0.243$ against $-0.191$. On the ordinary night it is zero ($0.000$ against $-0.023$). In summer, if one end of a link is connected mostly to well-connected zones, the other end tends to be connected mostly to poorly connected ones, more than the degrees require: the circles of the senders and those of the receivers are kept apart.

The first answer is therefore structural, and it is mostly about the beach. In summer the network turns into a network of senders and receivers, and most of its new mixing follows from that. None of this says who goes where. That is the next question.

\subsection{The roots pull every night, and in August the exodus generation pulls inland}\label{sec:res-roots}

On every night, ordinary or holiday, a city sends more of its residents to the provinces where more of them were born. Comparing two cities of the same province against the same destination, the pull of the roots, $\beta$ in equation~\eqref{eq:ppml}, is 0.53 on the ordinary night and 0.63 on 14--15 August in the 137 large mainland cities of the full census table (Section~\ref{sec:method}): a city with twice the share of residents born in a province sends 1.44 times as many people there on a weekday in February, and 1.55 times in August. The pull is significant on every night tested ($p\le0.002$).

Against the ordinary night, the summer increase of $+0.10$ is not secure ($p=0.18$). On this measure alone, the answer to the question would be ``a little, if at all'': holidays hardly strengthen the roots once the beaches take their share. The average hides two opposite movements, which the census groups separate (Figure~\ref{fig:roots}).

\begin{figure}[t]
\centering
\includegraphics[width=\linewidth]{fig3_roots_generation.pdf}
\caption{\textbf{The extra summer pull of home provinces comes from the exodus generation and goes inland.} Change in the pull of the roots ($\beta$ in equation~\ref{eq:ppml}) between 14--15 August and 13--14 February for each census group of 675 cities, with 95\,\% wild cluster bootstrap intervals (44 provinces); $+0.1$ is about 7\,\% more travellers per doubling of the group. \textbf{a}, Stays in inland and in coastal zones. \textbf{b}, Groups split by age. \textbf{c}, Each of the 117 pairs of one summer night and one February night.}
\label{fig:roots}
\end{figure}

Residents who arrived in their current province before 1981, most of them the generation of the exodus, pull harder in August: $+0.084$ (95\,\% interval 0.046 to 0.117, $p<0.001$). So do residents born in a village of 2,000 inhabitants or fewer: $+0.086$ (0.062 to 0.110). Residents who arrived in 2016--2023 pull less: $-0.072$ ($-0.114$ to $-0.020$). The extra pull goes only inland. For stays in inland zones, the increase is $+0.154$ for the pre-1981 group and $+0.123$ for the village-born; for stays on the coast it is $-0.017$ and $+0.019$, and neither can be told apart from no change. A doubling of the pre-1981 group born in a province brings about 11\,\% more of the city's residents to the inland part of that province in August than in February.

Three checks make this hard to explain away. The first is age: the exodus generation is old, and older people may simply travel differently. Splitting the groups by age, residents aged 65 or more who arrived after 1981 show no extra pull ($-0.006$), while those who arrived before 1981 show $+0.087$. What matters is when they came, not how old they are. The second is the choice of nights: the sign is the same in all 117 pairs of a summer night and a February night, and the size changes little when any one of the 44 provinces is left out ($+0.072$ to $+0.089$ for the pre-1981 group). The third is the baseline: against October weekdays instead of February, the increases are $+0.086$ and $+0.095$.

The other holidays point the same way with smaller sizes. On Good Friday the pre-1981 group adds $+0.057$ and the village-born $+0.068$. On All Saints and Christmas Eve the increases are smaller ($+0.026$ to $+0.039$), and at Christmas recent migrants no longer pull less.

\subsection{The emptied villages switch from senders to receivers}\label{sec:res-villages}

The villages that lost most people show the same pattern from the other side (Figure~\ref{fig:villages}). On an ordinary night, 24\,\% of the fifth of zones that lost most population receive more people than they send. On 14--15 August, 99\,\% do, and on Christmas Eve more than 99\,\%. Across all zones, the share of net receivers goes from 41\,\% to 65\,\% in August and 81\,\% at Christmas, while the large cities empty: Madrid's net balance goes from $+0.07$ to $-0.80$ in August.

\begin{figure}[t]
\centering
\includegraphics[width=\linewidth]{fig4_villages.pdf}
\caption{\textbf{Emptied villages switch from sending people to receiving them.} \textbf{a}, Net balance of each zone, (in $-$ out)/(in $+$ out) people sleeping away from home; white zones send more people than they receive. \textbf{b}, Share of zones that receive more than they send, by fifths of population lost since the 1900--1960 peak. Greys distinguish nights; in other figures colours mark coast, inland and cities.}
\label{fig:villages}
\end{figure}

The switch is not a matter of size, grouping or location. Holding fixed population, the number of municipalities in a zone, area, the coast, distance to the nearest city and province, the most emptied zones are 0.39 more likely than the least emptied to be net receivers on 14--15 August (95\,\% interval 0.23 to 0.55), and 0.44 more likely among zones made of a single municipality. On the ordinary night the difference goes the other way ($-0.12$, not significant). The summer result holds on all 9 sister nights and in 2022 and 2023.

The two holidays reach the same result in different ways. In August, the emptied zones receive about as many visitors as comparable zones that lost little population; what differs is that their own residents stay at home, with about half as many of them sleeping away. On Christmas Eve, visitors per resident rise as well. The stock of second homes fits this difference: it accounts for about half of the link between emptying and the summer rise in visitors (48\,\%, interval 22 to 82\,\%), and for almost none of it on Christmas Eve (2\,\%). This fits a simple reading, that in summer many of those who come back stay in a house of their own and at Christmas with family, but the data do not show where each person sleeps.

\subsection{What did not survive}\label{sec:res-failed}

Two results that looked strong in the first analyses did not pass the checks, and I report them because they changed the story.

The first was that emptied villages move into the core of the network on holiday: their $k$-core rose about fourfold in summer. Comparing zones of similar size and location, the opposite holds: the coreness of emptied zones falls behind (by 9.4 shells in August), and degree-preserving rewiring reproduces that fall almost exactly ($-8.9$). The $k$-core of a zone here follows from its degree, not from where it sits in the network.

The second was that cities and the zones their residents visit form ``summer territories'', communities that are not contiguous on the map, like archipelagos of a diaspora. Communities of the bipartite network of cities and destinations do look like that. But a null model that keeps each city's and each destination's total and the decay of flows with distance produces at least as many non-contiguous communities as the real network, on every night. The territories are what distance and size produce. Holidays change who goes where, as Section~\ref{sec:res-roots} shows, but not in a way that a community algorithm can separate from geography.
